\section*{Bab \ref{ch:g12:buffers-predominance} --- Larutan Penyangga dan Diagram Predominansi}

\begin{solution}{exo:g12:buffers-predominance:1}
\omterm{def:g9:acids-bases-ph:ph}{pH} 2: \ce{CH3COOH}. \omterm{def:g9:acids-bases-ph:ph}{pH} 4.76: keduanya dalam jumlah yang sama. \omterm{def:g9:acids-bases-ph:ph}{pH} 7:
\ce{CH3COO-}.
\end{solution}

\begin{solution}{exo:g12:buffers-predominance:2}
\omterm{def:g9:acids-bases-ph:ph}{pH} 4: $1/(1 + 10^{0.76}) = 0.15$. \omterm{def:g9:acids-bases-ph:ph}{pH} 6: $1/(1 + 10^{-1.24}) = 0.95$.
\end{solution}

\begin{solution}{exo:g12:buffers-predominance:3}
Kuning; hijau (di dalam trayeknya); biru.
\end{solution}

\begin{solution}{exo:g12:buffers-predominance:4}
Pada \omterm{def:g9:acids-bases-ph:ph}{pH} 7.4, di bawah 9.26: \ce{NH4+}. Pada \omterm{def:g9:acids-bases-ph:ph}{pH} 11: \ce{NH3}.
\end{solution}

\begin{solution}{exo:g12:buffers-predominance:5}
Indikator itu sendiri adalah \omterm{def:g12:ka-and-pka:bronsted}{pasangan asam--basa}: dalam jumlah besar,
indikator akan mengubah \omterm{def:g9:acids-bases-ph:ph}{pH} yang seharusnya ditunjukkannya.
\end{solution}

\begin{solution}{exo:g12:buffers-predominance:6}
$10^{\mathrm{pH} - 4.76}$: 0.1; 1; 10.
\end{solution}

\begin{solution}{exo:g12:buffers-predominance:7}
Fenolftalein (8.0--10.0) atau biru timol (8.0--9.1): trayeknya mencakup
8.7. Merah metil berubah antara 4.2 dan 6.3, jauh dari 8.7: warnanya
sudah berubah jauh sebelumnya.
\end{solution}

\begin{solution}{exo:g12:buffers-predominance:8}
$\mathrm{pH} = 4.76 + \log(0.10/0.20) = 4.76 - 0.30 = 4.46$.
\end{solution}

\begin{solution}{exo:g12:buffers-predominance:9}
\omterm{def:g9:acids-bases-ph:ph}{pH} 1: \omterm{def:g9:ions:ion}{kation} \ce{H3N+-CH2-COOH}, muatan $+1$. \omterm{def:g9:acids-bases-ph:ph}{pH} 6: \omterm{def:g9:ions:ion}{ion} zwitter, muatan
keseluruhan 0. \omterm{def:g9:acids-bases-ph:ph}{pH} 12: \omterm{def:g9:ions:ion}{anion} \ce{H2N-CH2-COO-}, muatan $-1$.
\end{solution}

\begin{solution}{exo:g12:buffers-predominance:10}
Air: dari 7.0 menjadi $-\log(\num{1.0e-4}/0.101) = 3.0$, turun 4 satuan.
Penyangga: dari 4.76 menjadi $4.76 + \log(9.9/10.1) = 4.75$, turun 0.01.
\end{solution}

\begin{solution}{exo:g12:buffers-predominance:11}
\ce{NH4+}/\ce{NH3} ($pK_a = 9.26$). Perbandingan
$[\ce{NH3}]/[\ce{NH4+}] = 10^{9.5 - 9.26} = 1.7$.
\end{solution}

\begin{solution}{exo:g12:buffers-predominance:12}
\omterm{def:g9:acids-bases-ph:ph}{pH} turun satu satuan ketika perbandingannya mencapai $1/10$:
$(10 - n)/(10 + n) = 0.1$, jadi $n = \qty{8.2}{mmol}$. Penyangga yang
mengandung lebih banyak setiap bentuk dapat meredam lebih banyak asam
atau basa untuk perubahan perbandingan yang sama: kapasitasnya lebih
besar.
\end{solution}

\begin{solution}{exo:g12:buffers-predominance:13}
Perbandingan $10^{5.0 - 4.76} = 1.74$; asamnya berjumlah
\qty{10}{mmol}, jadi \qty{17.4}{mmol} etanoat, yaitu \qty{174}{mL}
larutan itu.
\end{solution}

\begin{solution}{exo:g12:buffers-predominance:14}
Biru timol memiliki dua \omterm{def:g12:ka-and-pka:bronsted}{pasangan asam--basa}, dengan dua $pK_a$, sehingga
dua perubahan warna: merah pada \omterm{def:g9:acids-bases-ph:ph}{pH} 1, kuning pada \omterm{def:g9:acids-bases-ph:ph}{pH} 5, biru pada \omterm{def:g9:acids-bases-ph:ph}{pH} 10.
\end{solution}

\begin{solution}{exo:g12:buffers-predominance:15}
\omterm{def:g10:concentration-and-dilution:dilution}{Pengenceran} membagi kedua konsentrasi dengan 10: perbandingannya, dan
karena itu \omterm{def:g9:acids-bases-ph:ph}{pH-nya} (4.76), tidak berubah. \omterm{def:g12:ka-and-pka:strong-acid}{Asam kuat} yang diencerkan
sepuluh kali naik satu satuan \omterm{def:g9:acids-bases-ph:ph}{pH}.
\end{solution}

\begin{solution}{pb:g12:buffers-predominance:1}
\textbf{1.} \ce{CO2 + 2H2O <=> HCO3- + H3O+}.

\textbf{2.} Karbon dioksida terlarut (bersama air) adalah asamnya,
hidrogenkarbonat basanya.

\textbf{3.} Nilai tabel berlaku pada \qty{25}{\celsius} dalam air murni;
darah bersuhu \qty{37}{\celsius} dan penuh dengan \omterm{def:g9:ions:ion}{ion-ion} lain, yang
mengubah tetapannya (dan para ahli fisiologi menghitung seluruh karbon
dioksida terlarut).

\textbf{4.} Pada \qty{24}{mmol/L}, $0.024 \times 5.0 = \qty{0.12}{mol}$.

\textbf{5.} Sumbu \omterm{def:g9:acids-bases-ph:ph}{pH} yang dipotong pada 6.1: \ce{CO2} di bawahnya,
\ce{HCO3-} di atasnya.

\textbf{6.} Pada 7.4, di atas 6.1: hidrogenkarbonat.

\textbf{7.} Sedikit basa: \omterm{def:g9:acids-bases-ph:ph}{pH-nya} di atas 7.

\textbf{8.} $10^{7.4 - 6.1} = 10^{1.3} = 20$.

\textbf{9.} $24 / 20 = \qty{1.2}{mmol/L}$.

\textbf{10.} $10^{1.28} = 19$ dan $10^{1.32} = 21$.

\textbf{11.} $20/21 = \qty{95}{\percent}$.

\textbf{12.} $10^{7.6 - 6.1} = 10^{1.5} = 32$.

\textbf{13.} \omterm{def:g9:ions:ion}{Ion} hidrogenkarbonat:
\ce{HCO3- + H3O+ -> CO2 + 2H2O}.

\textbf{14.} $10^{7.1 - 6.1} = 10$.

\textbf{15.} Dari $24/20 = 1.2$ menjadi $24/10 = \qty{2.4}{mmol/L}$:
dua kali lipat. Paru-paru bernapas lebih cepat dan lebih dalam, untuk
mengembuskan kelebihan karbon dioksida.

\textbf{16.} \omterm{def:g12:ka-and-pka:strong-acid}{Asam lemah} dan basanya masing-masing dapat meredam apa yang
ditambahkan, dan \omterm{def:g9:acids-bases-ph:ph}{pH} hanya bergantung pada perbandingan keduanya; \omterm{def:g12:ka-and-pka:strong-acid}{asam
kuat} hanya akan mendorong \omterm{def:g9:acids-bases-ph:ph}{pH} ke satu arah.

\textbf{17.} Sama sekali tidak berubah: perbandingannya, dan karena itu
\omterm{def:g9:acids-bases-ph:ph}{pH-nya}, tetap sama.

\textbf{18.} Sekitar 20.
\end{solution}
